\section{Two finite adaptive steps}\label{gex:adaptive-section}

The first-stage boundary can itself be used in the plain moment.
We now make two such uses, each with a strict endpoint margin.
The arithmetic construction and the recovery source remain those of
\cite{OAI26}; the extension of the plain moment is proved in
\cite[Theorem 1.1]{GaoAdaptive26}.

\subsection{The extended moment and its physical consumer}

\begin{lemma}[Extended plain moment]\label{gex:adaptive-plain}
Retain all hypotheses and notation of \cite[Lemma 18.1]{OAI26}.
For \(13/18\le\kappa\le3/4\), assume the global family bound
\(\beta_*\le(1+\kappa)/2\). Let the element rows satisfy
\(0<q_k\ll Z^m\), with \(\psi_k(n)=\tau(n)\chi_n(k)\), and put
\(M=m+q\), where \(q\) bounds the displayed moving conductor union
before cancellation. For positive total prime length \(z\), assume
\(n_1+n_2+6\kappa z\le M\) and restrict to inducing characters
outside the prescribed finite group \(\Theta\). Then
\[
 \sum_{k\in\mathcal R_z}
 \left|S_{\psi_k}(n_1;W_1)S_{\psi_k}(n_2;W_2)
                  \prod_i Q_{\psi_k,i}\right|^2
 \ll Z^{M+\varepsilon}.
\]
At \(z=0\), bounded plain lengths are unrestricted and the rows remain
nonprincipal. The slot mesh precedes the fixed slot count, uniformly
on bounded length ranges. Both plains and every slot use the same
puncture, fixed within the row sum. All defining-modulus zeros,
including canceled and six-divisible factors, remain. Prime supports
are disjoint before masking, and their coefficients are fixed finite
combinations from \(\Theta\), independent of the varying row.
The original label ranges, finite seminorms and polynomial separated-height
costs are retained.
\end{lemma}

The proof in \cite{GaoAdaptive26} follows every recursive child of
\cite[Section 18]{OAI26}. The global input licenses the prime contour
\((1+\kappa)/2+e_{\rm an}\) for each nonprincipal child; it is not a
bound inferred from its parent's detector bin. The changed terminal
length is \(z\le3M/13\). The comparison has admission reserve
\(M/15\); its clipping errors are \(8e_{\rm cl}/5\) and
\(13e_{\rm cl}/5\), so choosing
\(\xi+13e_{\rm cl}/5<M/15\) preserves that reserve. Natural zeros,
the common Fourier measure, exceptional equal-product cancellations
and the backward seminorm allocation are unchanged.

For the physical rows the twist is fixed-ray, so there is no additional
moving radical in \(M\). Write \(\delta=2a-1\), \(q=\delta x\),
\(0\le x\le1/2\), and \(\alpha=5/6\). The actual detector witnesses obey
\[
 r+m\ge t-O(\varepsilon),\qquad r\le t+O(1/\log U),
 \qquad 1\le t\le3/2.
\]
For \(r<1,m<1/2\), the inverse and doubled-plain capacities are
\((1-r)/2\) and \((1-2m)/(6\kappa)\). Put
\[
 c=\frac1{3\kappa},\quad B=2-2cx,\quad
 D=B+1-x,\quad P=B(1-x),\quad
 \mathcal J=(\alpha-\delta)D+\delta P.
\]
Their count exponents are
\(1-\delta\{x+(1-x)r\}\) and
\(1-\delta\{cx+B(t-r)\}\). Balancing them and the long-inverse
branch gives
\begin{equation}\label{gex:adaptive-count}
\begin{gathered}
 R_{\rm short}(t)=1-\delta+\delta P(3/2-t)/D,\qquad
 L(t)=1-\delta+(\alpha-\delta)(t-1),\\
 t_*=1+\frac{\delta P}{2\mathcal J},\qquad
 R_*=1-\delta+\frac{(\alpha-\delta)\delta P}{2\mathcal J}.
\end{gathered}
\end{equation}
These formulas use both witness inequalities, not an equality for their
lengths. Uniformly \(D\ge2\), \(\mathcal J\ge25/39\), and the inverse
crossing has \(r_*\ge23/37\). Thus the existing inverse prime supply
\(7/37\) suffices; the plain supply is at most \(c/6\le1/13\).
Zero-capacity neighborhoods use the unweighted estimates.
Whole-product conjugation preserves the physical coefficients and zeros.
Rowwise selected real parts and heights are recovered by the same common
parameter-Sobolev measure, not by row-dependent prime coefficients
\cite[Proposition 4.1]{GaoAdaptive26}.

\subsection{Two exact endpoint certificates}

Set \(\theta_0=20999/24000\), supplied by
Proposition~\ref{gex:first-stage}. For each of the following stages use
\(\kappa=2\theta_{i-1}-1\), and the geometry
\[
 \ell_i=1/6+p_i,\quad h_i=13/16+3p_i/2,\quad
 l_{x,i}=17/48-p_i/2,\quad l_{y,i}=23/48-p_i/2.
\]
Here \(\ell_i\) is the total prime length for stage \(i\), not an
individual slot length. The latter are fixed after the moment mesh.
\begin{center}
\small
\begin{tabular}{@{}cccccc@{}}
\toprule
Stage & \(\kappa\) & \(c\) & \(p_i\) & \(m_i\) & \(\theta_i=7/8-p_i/4\)\\
\midrule
1 & \(8999/12000\) & \(4000/8999\) & \(1/5836\)
  & \(10^{-8}\) & \(20425/23344\)\\
2 & \(8753/11672\) & \(11672/26259\) & \(107121/625000000\)
  & \(10^{-9}\) & \(2187392879/2500000000\)\\
\bottomrule
\end{tabular}
\end{center}
Both \(\kappa\)'s lie in \([13/18,3/4]\). At stage 2 the prime contour
is \(\theta_1+e_{\rm an}\), licensed by the already established stage 1,
not by the output \(\theta_2\).

The principal exponent remains \(C(s)=s-11/16\). The literal central
endpoint, with the count \eqref{gex:adaptive-count}, is
\[
 E_i=-1/48+(2/3)\delta+\delta x/6-(13/16)(1-R_*)
       +p_i\{-1/4+\delta(1+x)+(3/2)R_*\}.
\]
We give a compact rational certificate on the full rectangle
\(0\le\delta\le5/6,\ 0\le x\le1/2\). Set \(y=1/2-x\) and
\[
\begin{gathered}
 D=5/2-c+(1+2c)y,\qquad P=1-c/2+2y+2cy^2,\\
 k_i=1/48-5p_i/4-m_i,\qquad v_i=1/16+y/6+p_i y.
\end{gathered}
\]
Then exact multiplication gives
\begin{equation}\label{gex:adaptive-certificate}
\begin{aligned}
 2\mathcal J(-E_i-m_i)&=A_i\delta^2+B_i\delta+C_i,\\
 A_i&=2(P-D)v_i+h_iP,\\
 B_i&=2\alpha Dv_i+2(P-D)k_i-h_i\alpha P,\qquad
 C_i=2\alpha Dk_i.
\end{aligned}
\end{equation}
For both rows of the table, rational expansion of these displayed
formulas gives \([y^j]A_i\ge1/4\) for \(0\le j\le3\), and
\[
 [y^j](4A_iC_i-B_i^2)>d_j\quad(0\le j\le4),\qquad
 (d_0,\ldots,d_4)=
 (10^{-8},1/25,19/100,13/50,3/40).
\]
All bounds are rational coefficient inequalities. In particular
\(A_i>0\) and \(4A_iC_i-B_i^2>0\) for \(y\ge0\). The identity
\(4A_i(A_i\delta^2+B_i\delta+C_i)
 =(2A_i\delta+B_i)^2+4A_iC_i-B_i^2\)
proves \(E_i<-m_i\). Stage 2 therefore retains the full
\(10^{-9}\) margin; no sampling or decimal comparison is used.

\subsection{Activation at the new scales}

For each stage repeat the same compensated probe
\eqref{gex:compensated-source}, including every marked subset, the
original disjoint windows and the full quotient-free correction.
The rescaled minima are
\(M'\ge1/2-3p_i>0\), \(l_{x,i}-d\ge3/16-3p_i/2>0\), and
\(l_{y,i}-d-11b/6\ge1/12-3p_i/2>0\).
Lemma~\ref{gex:completed-low}
and the physical Gram recovery apply since \(\ell_i<1/5\);
the tuple payment remains
\[
 -d+(5\ell_i-1+d)_+/8\le-7d/8.
\]
Thus the complete low bound is
\(Z^{3/16-p_i/4+\varepsilon}=Z^{C(\theta_i)+\varepsilon}\),
with the same finite profile, cusp, mask and height costs.

For moderate rows use the unchanged \(D_1(10e)\) contours and
\eqref{gex:adaptive-count}. The floor, terminal and intermediate
branches retain the bounds
\[
 -7/1200+15p_i/8,\qquad
 -1/48-\delta/16+15p_i/8,\qquad
 -49/14400+139p_i/150,
\]
respectively; each has a margin larger than \(m_i\).
The frequency slope is between \(4/25\) and \(2\).
Choose \(0<\zeta<5\ell_i-h_i=1/48+7p_i/2\); then
\(\ell_i/d>1/5>7/37\) for \(1/2\le d\le h_i+\zeta\).
Thus the fixed disjoint physical slots provide the required capacities.
Allocate the aggregate losses in \eqref{gex:central-loss-budget}
below \(m_i/8\): choose the mesh before the even fixed slot count,
then the allowances involving that count. The strict denominator and
selection margins make this allocation uniform in the targets.

The second Euler box is now
\((\Rea s,\Rea w,\Rea z)\ge(\theta_i,19/20,33/200)\).
With \(a_0=p_i/4\), the good- and row-prime defects in
\eqref{gex:local-euler-defect} are
\(q_p^{-363/200+a_0}\) and \(q_p^{-33/40+a_0}\);
the denominator decay is \(56/25-3p_i/2>2\).
These give normal convergence and divisor-product control as before.
The selected-factor error exponents
\(-\theta_i,-99/100,-273/200+5p_i/4,-47/50\)
are all at most \(-4/5\). Hence the same principal correction,
nonvanishing \(H_\eta\), and logarithmic normalizer are retained.
Its residue reserves are
\(l_{y,i}/20,h_i/600,(2/5)\min_j\ell_{i,j}>0\).
Normalization losses are reserved after the slots are fixed.

Small rows use \((\Rea s,\Rea w,\Rea z)=(7/8+8e,1/2,17/50)\),
which lies in \(D_1(3/8)\) and is right of every target zero by the
input theorem. Under \(\beta_*>\theta_i\), its shift from
\(\beta_*+e\) costs less than \(p_i/4+7e\). Together with the scale
cost \(51p_i/100\), the saving is
\(63/800-19p_i/25>1/16\), including bad-place and nontrivial-unit
rows. No use of \(\beta_*+e\) in \(D_1(3/8)\) below \(7/8\) is made.
For large rows, \(B_{0,i}=53/32-3p_i/4\). After fixing \(\zeta>0\),
choose a fixed absolute contour order \(v\) so that
\(B_{0,i}+(h_i+\zeta)(1+\varepsilon)-\zeta v+\varepsilon\)
is below the required exponent. These estimates retain all axes,
joins and rescaled profiles. For each target the internal finite
orders precede its height exponent \(\tau_\eta\), which precedes
the later external tail order. Constants and thresholds may depend
on the target; the remaining power reserve is common.

\subsection{Finite feedback and the full family}

We finish the proof of Theorem~\ref{gex:main}. At stage \(i\), assume
\(\theta_i<\beta_*\le\theta_{i-1}\), put
\(\Gamma_i=\beta_*-\theta_i\), and choose the low allowance
\(\omega=\Gamma_i/2\). All slots and power reserves precede the
individual target. The recovery of Section~\ref{gex:recovery-section},
with the foregoing parameters, produces the same signal
\eqref{gex:common-signal} and a common \(\sigma_i>0\) such that
\[
 |J_\eta(Z)|\ll_\eta Z^{C(\theta_i)+\omega},\qquad
 |J_\eta(Z)-f_\eta(Z)|\ll_\eta Z^{C(\beta_*)-\sigma_i}.
\]
The uniform logarithmic control used in this recovery comes from the
input's all-height zero location and retains its conductor and height
logarithms. It does not require uniform target-dependent thresholds.
Let \(\epsilon_i=\min(\Gamma_i/2,\sigma_i)>0\).
The Mellin transform of \(f_\eta\) continues
\(e^{(s-5/6)^2}H_\eta(s)/L_F^S(s,\eta)\) holomorphically to
\(\Rea s>\beta_*-\epsilon_i\), exactly as in the previous section.
The numerator is nonzero there. A primitive target zero sufficiently
near the supremum contradicts this continuation; attainment of the
supremum is unnecessary. This proves \(\beta_*\le\theta_i\).
Starting with Proposition~\ref{gex:first-stage} proves stage 1 and
then stage 2, without a contour at an unproved output boundary.

Finite induced Euler factors are nonzero for \(\Rea s>0\), so the
same result holds for every finite-order Hecke presentation.
For a Dirichlet character \(\chi\), its norm lift
\(\eta_\chi(\mathfrak a)=\chi(N\mathfrak a)\) satisfies, after deleting
the modulus support and \(3\),
\[
 L_F^{S_F}(s,\eta_\chi)
 =L^{S_{\Q}}(s,\chi)L^{S_{\Q}}(s,\chi\chi_{-3}).
\]
Both factors are holomorphic on \(\theta_2<\Rea s<1\); their nonzero
product and finite restoration give the Dirichlet assertion there.
The coarse input covers \(\Rea s\ge1\), with principal poles excepted.
There is no height restriction. Functional equations and closure under
conjugation reflect critical-strip zeros, yielding the lower boundary
\(1-\theta_2=312607121/2500000000\). This proves the full theorem.
The two strict gains are \(41/35016000\) and
\(39539/3647500000000\); only these two finite steps are used.
